% SOURCE docs/release/2026-10-29/source/docs/writing-repair/2026-10-29/src/7cc06a3bd2696253.tex CHAR 0
\chapter{Changing policies, censoring and events that end a decision}
\label{r3:policy-targets}
The history g-formula identifies a specified intervention, not every
quantity labeled a policy effect. A bank may increase an existing offer
probability, modify an employee's proposed credit limit, or stop making
decisions after repayment. These operations use different information
and produce different statistical functionals. Their distinctions also
determine which nuisance functions enter an influence function.

\section{Incremental interventions depend on the assignment law}
\label{r3:incremental}
For one binary decision write $e(x)=P(A=1\mid X=x)$ and
$m_a(x)=\E[Y\mid A=a,X=x]$. Multiplying the treatment odds by a
fixed $\delta>0$ gives
\[
 g_\delta(1\mid x)=\frac{\delta e(x)}{D_\delta(x)},\qquad
 D_\delta(x)=1+(\delta-1)e(x).
\]
This intervention retains impossible actions as impossible. At $e=0$
or $e=1$, it leaves the action unchanged. Under consistency and
exchangeability its mean is
\begin{equation}
 \psi_\delta(P)=\E[v_\delta(X)],\qquad
 v_\delta(x)=\frac{\delta e(x)m_1(x)+(1-e(x))m_0(x)}{D_\delta(x)}.
 \label{r3:incremental-target}
\end{equation}
The longitudinal construction replaces each propensity by this odds
transformation at its own history. Its identification under the stated
sequential assumptions is established in
\citet[Theorem 1 and Corollary 1]{kennedy2019}. The target is relative
to the population's behavior mechanism. Two banks with different $e$
can have different $\psi_\delta$ even when their outcome response
surfaces coincide.

For the unknown-propensity nonparametric model, direct differentiation
of Equation~\eqref{r3:incremental-target} gives
\begin{equation}
 \begin{split}
 \phi_\delta(O)={}&
 \frac{\delta A+1-A}{D_\delta(X)}\{Y-m_A(X)\}
       +v_\delta(X)-\psi_\delta\\
 &+\frac{\delta\{m_1(X)-m_0(X)\}}{D_\delta(X)^2}\{A-e(X)\}.
 \end{split}
 \label{r3:incremental-gradient}
\end{equation}
To verify the three components, factor a likelihood score into those
of $X$, $A\mid X$ and $Y\mid A,X$. Varying the last factor produces
the weighted outcome residual; varying the first produces the centered
conditional value. Since
$\partial v_\delta/\partial e=\delta(m_1-m_0)/D_\delta^2$,
varying the treatment factor produces the last term. Conditional mean
zero and this score identity verify the canonical gradient in the
nonparametric model with finite second moments. This reconstructs
\citet[Corollary 2]{kennedy2019}; the last term would be absent for
a supplied policy whose probabilities were held fixed as $P$ varied.
At $\delta=1$ the expression reduces to $Y-\E[Y]$, as it must:
the target is the observed mean. An arbitrary off-support $m_a$ can
be used at $e=0$ or $1$, since its coefficient vanishes almost surely.

For $\delta$ in a compact subset of $(0,\infty)$, the single-decision
likelihood ratio is bounded without imposing $e\in[\epsilon,1-\epsilon]$.
This avoids requesting unsupported actions; it does not make unsupported
fixed-treatment means identifiable. Products of the bounded ratios
can still grow exponentially with the number of decisions. Moreover,
the target itself depends on $e$: correct outcome regressions alone
with an inconsistent propensity do not generally recover it. The
source's second-order remainders and uniform inference conditions
must replace a generic appeal to double robustness
\citep[Sections 4.3--4.4 and Theorems 3--4]{kennedy2019}.

\section{Modifying a natural action is another operation}
\label{r3:modified-policy}
Let $A_t^{\mathrm{nat}}$ denote the action the assignment mechanism
would propose after the already modified past. A longitudinal modified
treatment policy (LMTP) applies
$A_t^d=d_t(A_t^{\mathrm{nat}},H_t^d)$. For example, reduce a proposed
limit by a fixed amount when the result remains within a supported
range. Observing the proposed action is part of the operational
information requirement; it cannot be replaced silently by a forecast
of what an employee might have proposed.

The distinction between this natural-action intervention and a fresh
stochastic draw with its induced conditional distribution matters.
In \citet[Assumptions 1--3 and Theorem 1]{lmtp2023}, standard
sequential randomization identifies the stochastic analogue. The
natural-action LMTP needs the stronger independence involving future
treatment innovations as well as future covariate innovations. In
the source's structural notation these conditions are respectively
\[
 U_{A,t}\perp U_{L,t+1:T}\mid H_t,\qquad
 U_{A,t}\perp (U_{L,t+1:T},U_{A,t+1:T})\mid H_t.
\]
The past is interpreted according to the source's sequential model.
The stronger restriction prevents an unrecorded cause of today's
natural action and future natural actions from invalidating the
extended intervention recursion. Both require that each transformed
action remain in the conditional observed support.

For a continuous action with density $e_t(a\mid h)$, suppose the
fixed map $d_t$ is invertible on pieces $I_j(h)$ with inverses $b_j$.
Its induced density is the pushforward
\begin{equation}
 e_t^d(a\mid h)=\sum_j
 \ind\{b_j(a,h)\in I_j(h)\}
 e_t(b_j(a,h)\mid h)\left|\frac{\partial b_j(a,h)}{\partial a}\right|.
 \label{r3:pushforward}
\end{equation}
Every inverse branch contributes. Omitting a branch where a modified
and an unmodified proposal coincide gives the wrong density ratio.
For discrete actions the corresponding expression is a sum of
$e_t(s\mid h)$ over all $s$ mapped to $a$. Piecewise smooth
invertibility, the transformed density and their role in regular
estimation are explicit in \citet[Assumption 4 and Theorem 2]{lmtp2023}.
A cap that maps an interval to one exact dose creates a point mass;
when the observed action is purely continuous, pointwise outcome
evaluation there is not a regular nonparametric mean-estimation problem.

Even at one time, substituting $e^d$ into a fixed-policy score misses a
term. For a fixed map $d$ and the identified functional
$\psi_d=\E[m(d(A,X),X)]$, change of variables gives the gradient
\begin{equation}
 \phi_d(O)=\frac{e^d(A\mid X)}{e(A\mid X)}\{Y-m(A,X)\}
                 +m(d(A,X),X)-\psi_d.
 \label{r3:lmtp-gradient}
\end{equation}
Outcome-law differentiation gives the first term. Differentiating the
joint law of the natural $(A,X)$ gives the second. Its action-dependent
part is exactly what disappears if the induced policy is incorrectly
held fixed. The source's longitudinal version replaces this second
term by a backward sequence of modified-action evaluations and
weighted residuals. If the map itself is learned from the analysis
sample, an additional distinction is necessary: conditional evaluation
of a frozen learned rule versus inference for the population optimization
functional. Neither Equation~\eqref{r3:lmtp-gradient} nor the
fixed-rule longitudinal score supplies that latter inference automatically.

\section{Censoring adds an observation intervention}
\label{r3:censoring}
Use $C_t=1$ to mean that the trajectory remains observed immediately
before action $t$, with $C_0=1$ and monotone loss thereafter. In this
section the clock is $H_t,A_t,C_{t+1},L_{t+1}$: the next measurement
is recorded only if $C_{t+1}=1$. Let
\[
 c_t(h,a)=P(C_{t+1}=1\mid C_t=1,H_t=h,A_t=a).
\]
The full-data outcome exists even if its record is lost. To identify
the mean under $g$ and complete observation, assume consistency,
sequential treatment exchangeability, and conditional independence of
retention from the future full-data trajectory given $(H_t,A_t,C_t=1)$.
Treatment and retention probabilities must be positive on histories
reached by this joint intervention. The covariate kernels in the
g-formula are then estimated among those retained at the corresponding
clock. This is the joint-treatment construction in
\citet[Section 21.5, printed pp.275--276]{hernan2020}.

Cancelling the action and retention factors of the recorded-data law
gives, for a fully observed terminal $Y$,
\begin{equation}
 \E[Y^{g,\bar C=1}]
 =\E\left[Y\prod_{t=0}^{T-1}
       \frac{g_t(A_t\mid H_t)}{e_t(A_t\mid H_t,C_t=1)}
       \frac{C_{t+1}}{c_t(H_t,A_t)}\right].
 \label{r3:censor-weights}
\end{equation}
For a path that is lost, define its summand as zero at its first loss;
unobserved later factors are not evaluated. Sequential cancellation is
also a proof that earlier covariates may be affected by treatment.
No independence of censoring marginally is required. If observation
is decided after a newly measured covariate, that covariate must instead
precede censoring in the factorization and enter its conditional model.
Using the wrong clock cannot be repaired by increasing model capacity.

Recursive augmentation follows by treating action and retention as
separate intervention nodes and applying the longitudinal gradient
to the resulting ordered kernels. A combined weight in an otherwise
unmodified regression is not a derivation of the required residuals.
The right-censored setting of \citet{sequentialdr2018} develops
sequentially doubly robust constructions for these models. Their
conditions concern the appropriate backward nuisance functions, not
arbitrary accuracy on the raw future outcome at every time.

For servicing, a portfolio transfer can be loss of observation if
repayment remains a meaningful endpoint. If keeping an account on the
bank's platform also changes its service and repayment, ``prevent
censoring'' has substantive effects. One must specify either an
observation-only monitoring intervention or that actual servicing
intervention. Administrative follow-up limits likewise restrict the
identifiable time horizon: weighting cannot reveal events beyond all
customers' maximum observed duration.

\section{Survival, competing risks and financial value}
\label{r3:competing-events}
Consider the first event after loan origination: default ($J=1$) or
full prepayment ($J=2$). Let $K$ be its discrete time and define,
under a policy $g$,
\[
 \lambda^g_{j,t}=P(K=t+1,J=j\mid K>t;g),\quad
 S_t^g=P(K>t;g),\quad S_0^g=1.
\]
The two next-period categories here are mutually exclusive, with a
specified tie convention if raw events share a timestamp. Multiplying
conditional survival probabilities and partitioning by first-event time
give
\begin{equation}
 S_k^g=\prod_{t=0}^{k-1}(1-\lambda^g_{1,t}-\lambda^g_{2,t}),\qquad
 F^g_1(k)=\sum_{t=0}^{k-1}S_t^g\lambda^g_{1,t}.
 \label{r3:competing-incidence}
\end{equation}
These are marginal policy hazards, after integrating histories among
policy-specific survivors. They can be obtained from the full
history g-formula or suitable weighted risk sets; inserting observational
hazards confounded by treatment is insufficient. If instead the
competing event is ordered first inside an interval, its hazard
$h_{2,t}$ and the subsequent conditional default hazard $h_{1,t}$
give $\lambda_{2,t}=h_{2,t}$ and
$\lambda_{1,t}=(1-h_{2,t})h_{1,t}$.
The ordered convention is used in
\citet[Section 2.2, equations (1)--(3)]{competing2026}; its 2026
preprint supplies additional estimation theory under its own clock
and assumptions, not evidence of banking performance.

The effect on default incidence differs from the effect on a
cause-specific hazard. In a two-period project example, a reference
policy has constant category probabilities $(.10,.70)$ for default
and prepayment. A new policy has $(.08,.02)$. Equation~\eqref{r3:competing-incidence}
gives default risks $.12$ and $.152$: the default hazard falls, but
the sharply lower prepayment rate leaves more loans exposed to later
default. This is not an estimation error. A default-only objective
also omits interest earned, recoveries, prepayment revenue and servicing
cost. The financial objective can instead be written
\[
 V(g)=\E^g\left[\sum_{t=0}^{T-1}\gamma^t\ind\{K>t\}R_t
                    +\gamma^K B_J\ind\{K\le T\}\right],
\]
where the reward timing and terminal event cash flow $B_J$ must avoid
double counting. Restricted mean event-free time on a grid of widths
$\Delta_t$ is $\sum_t\Delta_t S_t^g$ by integrating the survival
indicator. It measures time in the active state, not profit.

Treating prepayment as ordinary censoring would target a hypothetical
world without prepayment only under extra assumptions specifying how
that removal affects default. The observed competing-risk law alone
does not identify the latent default time after prepayment. In the
reference example, mechanically censoring prepayments and retaining
a $.10$ default hazard gives $1-.9^2=.19$, which is not the actual
default incidence $.12$. Hazard ratios also compare different
post-treatment survivor populations; they need not equal an average
individual causal effect even in a randomized experiment
\citep[Technical Point 8.1 and Section 17.2]{hernan2020}.

The next fitted examples return to a different endpoint: discounted
servicing reward. They show how the observed-history estimators behave
when that reward remains defined after loss of observation. They do not
replace the competing-event target above with a censoring target. The
eligibility discussion then returns to decisions whose feasible population
itself changes under the policy.



% SOURCE docs/release/2026-10-29/source/docs/writing-repair/2026-10-29/src/75d6d5a2bd09c7a2.tex CHAR 0
\section{A fitted observed-history route and its choice against a latent model}
\label{s4:fitted-history}
The longitudinal gradient becomes a practical alternative to a POMDP fit
when the recorded state is a sufficient observed history summary. Consider
the three-state servicing process of Section~\ref{r3:collections-study},
with independent customer trajectories and eight decisions per customer.
The state records repayment, arrears or severe distress; repayment is
absorbing and has only the zero action. The policy probabilities are
specified before sampling. The two policies' value difference is the
target, so the covariance between their estimated values must be retained.

Split whole trajectories into two folds. In the other fold estimate
action-specific transition tables by conditional counts and estimate
the behavior probabilities in each eligible state. The half-count
smoothing of sparse cells is fixed in advance. Compute initial $Q_t,V_t$
by backward recursion. On the evaluation fold form the cumulative policy
ratios and the discounted-reward analogue of Equation~\eqref{r3:long-score}.
Subtract the two policies' returns within each customer, concatenate the
folds and use the sample variance of these contrasts divided by the total
number of customers. Treating each customer-date row as independent would
discard precisely the covariance that this score represents.

For fixed horizon and a finite state/action space with positive occupation
probabilities and overlap on eligible states, table estimates have root-$N$
rates. The second-order longitudinal remainder is then $O_p(N^{-1})$.
Cross-fold evaluation and finite-dimensional consistency make the fitted
score difference negligible at root-$N$ scale. The iid trajectory CLT
therefore supplies the stated variance for fixed policies. Deterministic
actions in the absorbing state are compatible with this argument because
both evaluation policies take that same action there. The argument does
not cover growing histories or unobserved manager information.

\subsection{Bounded backward targeting as a competing estimator}
An alternative updates the continuation regression on each evaluation
fold while preserving its allowed financial range. Suppose the discounted
remaining reward lies in a known interval $[l_t,u_t]$. Put
$q_t=(Q_t-l_t)/(u_t-l_t)$, use its logit as an offset, and fit one
intercept $\epsilon_t$ by weighted logistic loss with cumulative weight
$W_t=\prod_{s\le t}g_s/e_s$. With the already updated $V_{t+1}$,
the score equation is
\[
 \sum_{i\in I}W_{it}\left[
 \frac{R_{it}+\gamma V_{t+1}^*(S_{i,t+1})-l_t}{u_t-l_t}
 -\operatorname{expit}\{\operatorname{logit}(q_t(S_{it},A_{it}))+\epsilon_t\}
 \right]=0.
\]
Map back to dollars and average the updated $Q_t^*$ over the supplied
policy to obtain $V_t^*$. This is a weighted-intercept fluctuation, not
the unweighted clever-covariate fluctuation described earlier. Both solve
the relevant weighted residual equation; they need not yield identical
finite-sample fits. Backward updating keeps the pseudo-outcome consistent
with the continuation function used at the next date. The bounds used here
include the worst possible action cost and terminal loss; they are not
learned from sample minima and maxima.

Summing the fitted score equations makes each fold's mean augmentation
term zero up to the root solver's tolerance. Thus its targeted plug-in and
augmented estimate agree. Evaluation-fold fitting here is one-dimensional
at each of a fixed number of dates, so finite-dimensional empirical-process
control preserves the first-order score expansion. The variance still uses
the full targeted influence values, not just the variance of initial-state
predictions. For flexible evaluation-fold regressions this argument would
need replacement; the fact that a targeting equation was solved is not
sufficient by itself.

In 500 experiments with 3,000 trajectories, changing action probabilities
from $(0,.2,.4)$ to $(0,.5,.8)$ has a true gain of $46.03$ per customer.
Augmentation and targeting both cover it in 96.8\% of runs, with RMSE
$4.47$ and $4.46$. A direct fitted recursion has RMSE $2.87$, benefiting
from the correctly specified small transition model, while cumulative
weighting has RMSE $4.83$. These latter comparisons concern point error;
their fitted-model/propensity uncertainty was not omitted from a purported
confidence interval. No interval is reported for them in this experiment.

Changing instead from never assisting to assisting every eligible state
has gain $137.24$. Augmentation covers in 95.6\% and targeting in 93.6\%
of runs, with RMSE $29.02$ and $24.39$. Direct recursion has RMSE $8.64$
and weighting $45.44$. The same sample size supports a modest stochastic
change much better than the deterministic policy contrast. All numerical
failures are counted; none occurred in these runs. The 500-replication
coverage standard errors are about one percentage point near 95\%.

The method hierarchy follows the information claim. If recorded history
is sufficient, an augmented or targeted observed-history route avoids
latent-state identification assumptions and allows robust nuisance fitting.
A small correct transition model can improve precision, but its advantage
in a simulated Markov population is conditional on that restriction.
If recorded history leaves important hidden state, an identified controlled
latent model can add information through measurement and dynamics, at the
cost of their rank, likelihood and assignment assumptions. In the binary
controlled experiment of Section~\ref{s4:latent-inference}, 100 randomized
replications give 96\% coverage of the $46.31$ gain; the Monte Carlo
standard error is $0.020$. Reducing emission separation gives 95\% coverage
over 60 replications, but repayment rewards remain state-informative, so
this is not evidence of uniform performance near singular full-model
information. An omitted private-state logger gives zero coverage and
bias $-161.26$ over 60 replications despite no numerical interval failure.
A successful optimizer cannot choose between these causal information
models on the bank's behalf.



% SOURCE docs/release/2026-10-29/source/docs/writing-repair/2026-10-29/src/ecee5558461e39fc.tex CHAR 0
\section{Continuous histories, dropout and separately learned policies}
\label{s5:history-censoring}
Finite tables are useful only when their state captures the information
that jointly determines action and outcome. A servicing system may instead
observe utilization, recent payment movements, earlier assistance and the
number of previous interventions. Treat the resulting history $H_t$ as
continuous. Let $D_t$ indicate that reward $R_t$ and next history $H_{t+1}$
are observed, conditional on remaining observed before decision $t$.
Censoring is monotone; default itself is an outcome, not dropout. The
target intervention replaces action probabilities by $\pi_t(a\mid h)$
and ensures follow-up. This joint intervention requires an argument that
follow-up does not itself change the customer process, or a model for how
it does. It is not the mean among customers who happened to remain visible.

Write $e_t(a\mid h)$ for the historical action probability and
$c_t(h,a)$ for conditional survival of observation. Consistency,
sequential exchangeability for actions and censoring, and positive
$e_tc_t$ on the policy's support identify
\begin{align}
 Q_t(h,a)&=\E[R_t+\gamma V_{t+1}(H_{t+1})\mid H_t=h,A_t=a,D_t=1],\\
 V_t(h)&=\sum_a\pi_t(a\mid h)Q_t(h,a),\qquad V_H=0,\qquad
 v_\pi=\E[V_0(H_0)].
 \label{s5:censor-recursion}
\end{align}
Exchangeability must concern the future histories used by later decisions,
as well as the terminal reward. Under the observed-data model the trajectory
score is
\begin{equation}
 U_\pi=V_0(H_0)+\sum_{t=0}^{H-1}\gamma^t W_t
       \{R_t+\gamma V_{t+1}(H_{t+1})-Q_t(H_t,A_t)\},\quad
 W_t=\prod_{s=0}^t\frac{\pi_s(A_s\mid H_s)D_s}
 {e_s(A_s\mid H_s)c_s(H_s,A_s)}.
 \label{s5:censor-score}
\end{equation}
Every summand after dropout is defined to be zero without accessing the
unobserved reward/history. The implemented check changes those hidden
values arbitrarily and verifies an identical score. This is a substantive
observation constraint, not merely a missing-value imputation convention.

\subsection{Telescoping, the remainder and the machine-learning obligation}
Conditional on the fitted nuisances, put
$\hat r_t=e_tc_t/(\hat e_t\hat c_t)$ and
$\delta_t=\hat Q_t-Q_t$, with
$\delta V_t=\sum_a\pi_t(a\mid H_t)\delta_t(H_t,a)$.
Changing measure one step at a time in the fitted score gives
\begin{align*}
 \E[\hat U_\pi]-v_\pi
 &=\E[\delta V_0]
   +\sum_t\gamma^t\E^{\pi}
       \left[\prod_{s\le t}\hat r_s
       \{\gamma\delta V_{t+1}-\delta_t\}\right]\\
 &=\sum_t\gamma^t\E^{\pi}
       \left[\prod_{s<t}\hat r_s(1-\hat r_t)\delta_t\right].
\end{align*}
Here $\E^\pi$ is the full-follow-up interventional path expectation; the
empty product is one. The second equality pairs the time-$t$ regression
error with the continuation error at time $t-1$, integrating the current
action under $\pi$. It exposes exactly why both the treatment and censoring
mechanisms belong in the weight. It also proves zero bias if every $Q_t$
is correct, or every product $e_tc_t$ is correct. More selective robustness
requires the appropriate backward consistency of the estimated continuation
regressions; it is not obtained by declaring an arbitrary collection of
models correct after fitting.

For fixed horizon, bounded policy ratios, and nuisance probabilities
bounded away from zero, Cauchy--Schwarz bounds this remainder by
\[
 C\sum_t\|\hat Q_t-Q_t\|_2
   \{\|\hat e_t-e_t\|_2+\|\hat c_t-c_t\|_2\}.
\]
The norms must dominate the relevant policy occupations; fixed overlap
provides the required change-of-measure bounds. Cross-fitting whole
customers controls the evaluation empirical process if the fitted score
converges in $L_2$. Product error $o_p(N^{-1/2})$, finite score variance and
an independent-trajectory CLT then give
$\sqrt N(\hat v_\pi-v_\pi)=N^{-1/2}\sum_i(U_{\pi,i}-v_\pi)+o_p(1)$.
For two policies use $U_{\pi_1,i}-U_{\pi_0,i}$ within each customer and its
variance. Summing two separately estimated variances would lose their
covariance. Fixed clipping is justified only if it is asymptotically inactive
under the asserted overlap; otherwise its bias or changed target must be
analyzed. A fitted probability with a good calibration plot does not prove
exchangeability, overlap in unobserved histories, or the product-rate bound.

The estimator is compatible with flexible regression, but its theorem is
conditional on those rates. Histogram boosting in the experiment is an
implemented nuisance choice, not a general proof of its rates for arbitrary
customer histories. Recent functional-history work makes a similar
distinction: a learned representation must preserve the outcome and mechanism
information needed by the score \citep[Definition 1, Assumptions 5--9,
Theorems 3--4]{ran2026}. Its approximation errors enter products alongside
estimation error. A representation with small prediction loss can still
discard a confounder. Its target concerns decisions based on observed
irregular measurements; a policy based on an unobserved complete trajectory
needs a further measurement model. The source's mean-aggregation Wald result
also does not automatically extend to its separate robust point aggregators.

\subsection{Selection, censoring and the fitted comparison}
An independent randomized pilot of 2,500 synthetic customers learns a
four-coefficient immediate-reward contrast. The resulting stochastic policy
uses action probability .7 where that fitted contrast exceeds .5 and .2
elsewhere; the reference policy always uses .2. This is a fixed learned
candidate after the pilot, not a proof of dynamic optimality. Conditional
on the pilot, the target is its actual policy value. An unconditional claim
about a policy-learning algorithm would require a different sampling
experiment that repeats the pilot as well. Regret relative to an optimal
policy is another estimand.

Evaluation uses 2,400 fresh customers, four decisions and discount .97.
Histories include a baseline covariate, continuous current and lagged risk,
previous action and cumulative action count. These summaries are sufficient
by the declared generator; no generic learned-state sufficiency is assumed.
The action and observation probabilities are fitted with logistic regressions;
arm-specific histogram boosting estimates continuation rewards backward.
Two whole-trajectory folds keep each evaluated customer's nuisance fits
outside that fold. A separate million-trajectory intervention simulation
gives contrast $3.5813$ with Monte Carlo standard error $.0040$, much smaller
than the evaluation interval standard error of about $.295$.

Across 120 repetitions under coarsening at random, the full score has bias
$-.0012$, RMSE $.2873$ and 95.8\% coverage at nominal 95\%. Ignoring
censoring gives bias $-.0199$ and 83.3\% coverage: the continuation model
is useful, yet its score variance is no longer an adequate inferential
construction. Direct recursion has mean bias $-.0376$; it is a point
challenger without a claimed interval. In the adverse generator, observation
also depends on the otherwise unseen reward innovation. The full fitted
score now has bias $-.1434$ and 89.2\% coverage; ignoring censoring covers
57.5\%. This is an observed failure of the identifying censoring condition,
not proof that a different learner will repair MNAR data. At 120 repetitions
the Monte Carlo standard error of coverage near 95\% is about .020.

Irregular opportunities require a further distinction. One may model the
timestamps and elapsed time as part of each observed history and retain the
natural visit clock; changing when visits occur is a different intervention.
The event-intensity likelihood and weights in Chapter~\ref{r3:continuous}
provide that second route. Interpolating missing months and treating them as
observed decisions supplies neither exchangeability nor the visit likelihood.
Thus the proposed servicing workflow first selects among observed-history
adjustment, an explicit controlled measurement model, and a proxy bridge.
Their identifying information differs even when all three output a policy
value in the same currency.


% SOURCE docs/release/2026-10-29/source/docs/writing-repair/2026-10-29/src/7cc06a3bd2696253.tex CHAR 13191


\section{Eligibility changes the population as well as support}
\label{r3:eligibility}
Let $E_t$ indicate eligibility before action $t$ and require
$g_t(\mathcal A(H_t)\mid H_t)=1$ for the feasible action set.
An absorbing repaid state can have one action, ``no further offer''.
Its assignment is deterministic and causes no positivity defect for
policies that also take that action there. Requesting assistance after
repayment is an infeasible intervention, not a rare-treatment problem.

An outcome such as satisfaction with an assistance offer may be
undefined when no offer is possible. A policy-specific eligible mean
can be expressed as a ratio
\[
 \theta_g=\frac{N_g}{D_g},\quad
 N_g=\E[E_t^gY_t^g],\quad D_g=P(E_t^g=1)>0,
\]
where the product denotes an explicitly defined zero contribution
outside eligibility. It does not impute satisfaction for an ineligible
person. If $N_g,D_g$ are identified differentiable functionals with
gradients $\phi_N,\phi_D$, the quotient rule yields
\[
 \phi_{\theta_g}=\frac{\phi_N-\theta_g\phi_D}{D_g}.
\]
Thus rare eligibility amplifies uncertainty even if treatment support
within the eligible subset is adequate. Comparing $\theta_g$ and
$\theta_{g'}$ compares different populations if the policies change
eligibility. The always-eligible principal stratum
$\{E_t^g=E_t^{g'}=1\}$ is a cross-world object and is not recovered
by separately estimating these two ratios.

The eligible treatment effect of \citet[Section 2.2, equation (1)]{selective2025}
holds the prior treatment path fixed and contrasts the current action
among units eligible under that same prior path. Its no-anticipation,
sequential-ignorability and overlap assumptions are stated on eligible
histories; they do not require randomizing eligibility itself. The
paper indexes sequential eligible opportunities with monotone dropout.
A bank's calendar process in which eligibility disappears and returns
needs its own state and event mapping before applying those formulae.

These distinctions specify the operations a longitudinal product must
accept. An action policy, a natural-action modification, an
observation intervention and an event-removal hypothetical are separate
scientific inputs. An interface that records only an outcome column,
an action column and a horizon cannot distinguish their estimands.

